\printedpage{9}
% head: A. Grothendieck : Dérivateurs, Ch. I, Édition des Universités de Montpellier II et Paris VII
\noindent suggests proceeding in the dual way). But the two coefficients one will find by these two procedures, products either of
\[
\begin{tikzcd}
F \arrow[dr] & \\
 & i_*i^*F \\
e \arrow[ur] &
\end{tikzcd}
\quad ,
\]
or of
\[
\begin{tikzcd}
 & F \\
j_!j^*F \arrow[ur] \arrow[dr] & \\
 & e
\end{tikzcd}
\quad ,
\]
should be \textit{canonically} isomorphic.

\noindent\msp{page 11}If $\DD$ is a derivator, the category
\[
\mathcal{A} = A_{\DD} = \DD(e) % edition reading verified: plain italic A_D is retained; this alphabet variation is not established to be an error
\]
plays an important role. It is called the \textit{base category,} or \textit{the fundamental category}, of the derivator, these objects are the \textit{fundamental coefficients.} For arbitrary $I$ in $\Dia$, we thus have for $F \in \DD(I)$, coefficient on $I$ % [sic] ces objets -> ses objets, possessive intended
\[
I^\circ \longrightarrow \mathcal{A} \quad , \qquad i \longmapsto F_i
\]
a diagram of type $I^\circ$ in $\mathcal{A}$, functorially, \textit{i.e.}
\[
\DD(I)\longrightarrow
\uHom(I^\circ,\underset{\substack{\Vert\\\DD(e)}}{\mathcal A})
\]
% the vertical equality sign is printed under the letter \mathcal{A} of \uHom(I^\circ,\mathcal{A}), with \DD(e) beneath it
(\textit{cf.} page 8).

\noindent\textbf{Der}$_3$ : For every $I \in \Dia$, this functor is \textit{conservative,}\footnote[4]{N.B. But this functor has no desire to be \textit{faithful.}} \textit{i.e.} for
\[
F \longrightarrow G
\]
\printedpage{10}%
% head: A. Grothendieck : Dérivateurs, Ch. I, Édition des Universités de Montpellier II et Paris VII
in $\DD(I)$ to be \texttt{[an]} isomorphism, it is necessary and sufficient that for every $i \in I$,
\[
F_i \longrightarrow G_i
\]
be one.

\noindent\textbf{Corollary.} \textit{The canonical functor (for $I$, $J \in \Dia$)}
\[
\DD(I \times J) \xrightarrow{\ \theta\ } \uHom(J^\circ, \DD(I))
\]
\textit{is conservative.}

Indeed, if $F$ is in $\DD(I \times J)$, we have for $(i,j) \in I \times J$
\[
F_{i,j} \simeq \theta(F)(j)_i \quad .
\]

\noindent\msp{page 12}Intuitively, one regards a derivator $\DD$ as ``being'' the fundamental category $\mathcal{A}$, with an additional ``structure'', \texttt{[to]} wit the datum of the ``extensions'' of $\mathcal{A}$ (as category of point coefficients) to the diagrams $I \in \Dia$, contravariantly in $I$. A coefficient $F$ in $\DD(I) \defeq \mathcal{A}_I$ is regarded as ``being'' the corresponding diagram $(F_i)_{i\in I}$ of fundamental objects, with an additional datum of ``globalization''. The $F_i$, or ``fibers'' of $F$ at the $i \in I$, are the local characteristics, of the ``global on $I$'' coefficient $F$. % [sic] caractéristiques locaux -> caractéristiques locales, agreement (feminine)

All known derivators are obtained in the following standard way. One starts from a ``category of models''
\[
\mathcal{M} \quad , \quad \Sigma
\]
equipped with a set of arrows, called ``equivalences'', and satisfying the relevant axioms for passing to a category of fractions
\[
\mathcal{A} = \mathcal{M}\Sigma^{-1}
\]
which will be the ``fundamental category''. One defines for every $I$ in $\Dia$
\[
\begin{aligned}
\mathcal{M}_I &= \uHom(I^\circ,\mathcal{M}) \\
\Sigma_I &= \{u \in \mathsf{Fl}(\mathcal{M}_I) \mid \forall i \in I, u_i \in \Sigma\} \\
\mathcal{A}_I &= \mathcal{M}_I\Sigma_I^{-1}
\end{aligned}
\]

\noindent\msp{page 13}It is clear that $\mathcal{A}_I \defeq \DD_{\mathcal{M},\Sigma}(I)$ depends on $I$ in a contravariant functorial way, and even in a 2-functorial way.

Of course, to really have a derivator, one must check the axioms. The strongest is the existence of the $f_!$, $f_*$. It is proved by Thomason when $(\mathcal{M},\Sigma)$ comes from a ``model category'' in Quillen's sense. It is not clear to me that then $(\mathcal{M}_I,\Sigma_I)$ also comes from a model category (in which for example the fibrations \printedpage{11}%
% head: A. Grothendieck : Dérivateurs, Ch. I, Édition des Universités de Montpellier II et Paris VII
and cofibrations would be defined from these data in $\mathcal{M}$, as earlier for the equivalences).

When one takes for $(\mathcal{M},\Sigma)$ a category of models for $\mathsf{Hot}$ (for example the semi-simplicial sets $\mathcal{M} = \uHom(\bDelta^\circ,\Ens)$, whence $\mathcal{M}_I \simeq \uHom(\bDelta^\circ \times I^\circ,\Ens) \simeq \uHom(\bDelta^\circ,\underbrace{\uHom(I^\circ,\Ens)}_{\widehat{I}})$, category of semi-simplicial objects in the topos $\widehat{I}$, one finds a description of a derivator (nonadditive), which will play a particularly important role. I denote it by % Resolved: exactly two closing parentheses in the formula; the source's outer prose parenthesis remains unclosed.
\[
\mathsf{HOT} \quad ,
\]
its category of fundamental coefficients is $\mathsf{Hot}$, the standard homotopy category.

\noindent\msp{page 14}The notion of ``morphism of Derivators'' is clear, it is a morphism between two strict 2-functors
\[
\Dia^\circ \longrightarrow \mathcal{C}at
\]
Thus
\[
u : \DD \longrightarrow \DD'
\]
is defined by giving, for every $I \in \Dia$,
\[
u_I : \DD(I) \longrightarrow \DD'(I)
\]
with commutativity (strict?) for $f : I \to J$
\[
\begin{tikzcd}
\DD(J) \arrow[r, "f^*_{\DD}"] \arrow[d, "u_J"'] & \DD(I) \arrow[d, "u_I"] \\
\DD'(J) \arrow[r, "f^*_{\DD'}"] & \DD'(I)
\end{tikzcd}
\quad . \tag*{(1)}
\]
We have
\[
\begin{tikzcd}[column sep=large]
\uHom(I,J)^\circ \arrow[r, "\DD_{I,J}"] \arrow[d, "\DD'_{I,J}"'] & \uHom(\DD(J),\DD(I)) \arrow[d, "\phi \mapsto u_I \circ \phi"] \\
\uHom(\DD'(J),\DD'(I)) \arrow[r, "\psi \mapsto \psi \circ u_J"'] & \uHom(\DD(J),\DD'(I))
\end{tikzcd}
\tag*{(2)}
\]
and the commutativity of (1) means the commutativity of (2) \textit{for the objects}. But of course we want it for the arrows too, \textit{i.e.} we want the commutativity of diagram (2).\footnote[5]{Give the example of a morphism of derivators defined by a functor of categories of models $(\mathcal{M},\Sigma) \to (\mathcal{M}',\Sigma')$.}

Probably we want a little more from morphisms of derivators - we want commutation with the $f_*$ and (?) with the $f_!$ (this time only up to canonical isomorphism) \ldots\footnote[6]{Commutation with $f_*$ (resp. $f_!$) \textit{left exact} (resp. \textit{right exact}) morphism. The \textit{exact} morphisms are those that are both left exact and right exact.}

\printedpage{12}
% head: A. Grothendieck : Dérivateurs, Ch. I, Édition des Universités de Montpellier II et Paris VII
\noindent\msp{page 15}For example, consider the derivator
\[
\mathsf{HOTAB}
\]
whose fundamental category is $\mathsf{Hotab}$, the full subcategory of the derived category $\mathsf{Der}(\mathsf{Ab})$ of $\mathsf{Ab}$, consisting of the complexes with vanishing cohomology in degree $> 0$. One can describe $\mathsf{HOTAB}$ by the category of models $\uHom(\bDelta^\circ,\mathsf{Ab})$ of semi-simplicial abelian groups, with as equivalences the quasi-isomorphisms in the usual sense. The canonical ``abelianization functor''
\[
\mathsf{Hot} \longrightarrow \mathsf{Hotab}
\]
extends to a morphism of derivators
\[
\mathsf{HOT} \longrightarrow \mathsf{HOTAB} \quad .
\]
The question of its compatibility with the $f_*$ does not seem tautological. It would imply, for example, that if
\[
\begin{tikzcd}
X \arrow[d, "f"'] & X' \arrow[l, "g"'] \arrow[d, "f'"] \\
Y & Y' \arrow[l, "j"]
\end{tikzcd}
\]
is a cartesian diagram in $\mathcal{M} = \uHom(\bDelta^\circ,\Ens)$ with $f$ a fibration (of Thom in the setting of topological spaces, with $f$ for example a locally trivial fibration) then the diagram in $\mathsf{Hotab}$
\[
\begin{tikzcd}
\mathbb{R}\Gamma_!(X) \arrow[d] & \mathbb{R}\Gamma_!(X') \arrow[l] \arrow[d] \\
\mathbb{R}\Gamma_!(Y) & \mathbb{R}\Gamma_!(Y') \arrow[l]
\end{tikzcd}
\]
is $\mathsf{hot}$-cartesian\footnote[7]{This is false already if $Y$ is the one-point space : trivially the $\mathbb{R}\Gamma_!$ does not transform products into sums! This functor is surely \texttt{[not]} ``left exact'' (as a morphism of derivators).}. And likewise for a $\mathsf{hot}$-amalgamated sum ($\mathsf{hot}$-cocartesian square in $\mathsf{Hot}$). % [sic] carrée -> carré, agreement (masculine noun)

\noindent\msp{page 16}If this is true, and not known to boot, it would be scandalous all the same!

Proceeding as for the base change morphisms for the $Rf_*$, we find in any case a canonical morphism, for $f : I \to J$
\[
u_Jf_*^{\DD} \longrightarrow f_*^{\DD'}u_I
\]
\printedpage{13}%
% head: A. Grothendieck : Dérivateurs, Ch. I, Édition des Universités de Montpellier II et Paris VII
and likewise
\[
f_!^{\DD'}u_I \longrightarrow u_Jf_!^{\DD} \quad ,
\]
and the question is whether these are isomorphisms. Analogous questions for the $f^!$, for $f$ a closed immersion (or for $f^?$, for $f$ an open immersion). Since the notion of a morphism of derivators ought all the same to ``fit'' the well-known notion of a (triangulated) morphism of triangulated categories\footnote[8]{No! The case of triangulated categories is ``misleading'', because for these left and right exactness amount to the same thing.}, hence be exact in a suitable sense, this does show that one must require a morphism of derivators to commute with the $f_*$ and with the $f_!$.\footnote[9]{N.B. One can conceive of a functor between derivators being left exact or right exact or both.} I should ask Thomason whether he has looked, for a morphism
\[
\mathcal{M} \longrightarrow \mathcal{M}'
\]
between categories of models, which is say \textit{left exact}, and ``compatible'' with the equivalences, fibrations (and cofibrations (???)).

\noindent\textsc{Notation} : $\mathsf{DER}$ for the 2-category of \textit{derivators.}

\noindent\msp{page 17}\textbf{Duality}. If $\DD$ is a derivator, we are going to define a \textit{dual} derivator $\DD^\circ$ by
\[
\DD^\circ(I) \defeq \DD(I^\circ)^\circ \quad .
\]
The coefficients on $I$ of type $\DD^\circ$ are the coefficients on the category $I^\circ$ opposite to $I$ of type $\DD$, and as a category it is the category \textit{opposite} to that of the coefficients of type $\DD$ on $I^\circ$. To define the 2-functorial law, we note that for $I$, $J$ in $\Dia$ we have % [sic] les coefficient -> les coefficients, plural agreement
\[
\begin{tikzcd}[column sep=large]
\uHom(I,J)^\circ \arrow[r, "\DD_{I,J}"] \arrow[d, "\wr"', "\text{duality}"] & \uHom(\DD(J),\DD(I)) \arrow[d, "\wr"', "\text{duality}"] \\
\uHom(I^\circ,J^\circ) \arrow[r, "\theta"] & \uHom(\DD(J)^\circ,\DD(I)^\circ)^\circ
\end{tikzcd}
\quad ,
\]
where $\theta$ is the unique functor making the diagram commutative. Then $\theta^\circ$is a functor
\[
\uHom(I^\circ,J^\circ)^\circ \longrightarrow \uHom(\DD(J)^\circ,\DD(I)^\circ) \quad ,
\]
that is by definition $\DD^\circ_{I,J}$ \texttt{[rather} $\DD^\circ_{I^\circ,J^\circ}$\texttt{]}.

If $\DD$ is defined by $(\mathcal{M},\Sigma)$, then $\DD^\circ$ by $(\mathcal{M}^\circ,\Sigma^\circ)$, since
\[
(\mathcal{M}^\circ)_I \defeq \uHom(I^\circ,\mathcal{M}^\circ) \simeq \uHom(I,\mathcal{M})^\circ \simeq (\mathcal{M}_{I^\circ})^\circ
\]
% foot (pages 9--13): Version 8 juillet 2001